\documentclass[12pt,a4paper]{article}
\usepackage[utf8]{inputenc}
\usepackage{xeCJK}
\usepackage{geometry}
\usepackage{booktabs}
\usepackage{array}
\usepackage{longtable}
\usepackage{multirow}
\usepackage{graphicx}
\usepackage{fancyhdr}
\usepackage{color}
\usepackage{xcolor}
\usepackage{amsmath}
\usepackage{amsfonts}
\usepackage{amssymb}

\usepackage{pdfpages}

\geometry{left=2cm,right=2cm,top=2.5cm,bottom=2.5cm}


\setCJKmainfont{Noto Serif CJK TC}
\setCJKsansfont{Noto Sans CJK TC}
\setCJKmonofont{Noto Sans CJK TC}


\pagestyle{fancy}
\fancyhf{}
\rhead{\thepage}
\lhead{心臟內科見習學習歷程成果報告}

\title{\Large \textbf{心臟內科見習學習歷程成果報告}\\
\large Clinical Shadowing Program Assessment Report}
\author{指導醫師：謝慕揚 醫師\\
見習學生：周煜珉 醫師}
\date{見習期間：2025年9月2日 - 2025年9月12日}

\begin{document}

\maketitle

\section{執行摘要}

周煜珉醫師於2025年9月2日至9月12日期間完成為期兩週的心臟內科臨床見習計畫。本報告基於學生的每日學習記錄、既定教學計畫以及心臟內科可信賴專業活動(EPAs)評估標準，對學習成果進行綜合評估。

\section{計畫執行與實際學習比較}

\subsection{計畫達成度分析}

\begin{table}[h]
\centering
\caption{教學計畫與實際執行對照表}
\begin{tabular}{|p{3cm}|p{5cm}|p{6cm}|}
\hline
\textbf{計畫項目} & \textbf{原定安排} & \textbf{實際執行情況} \\
\hline
門診見習 & 週一、週三門診觀摩 & 實際參與週一門診，觀摩70+病例，涵蓋多種疾病類型 \\
\hline
病房迴診 & 每日病房迴診與病患評估 & 持續參與病房迴診，接診多位新病患，完成病歷撰寫 \\
\hline
心導管室實習 & 觀摩基礎至進階手術 & 觀摩PCI、PTA、TPM、adrenal vein sampling等多種術式 \\
\hline
超音波訓練 & 標準切面學習與操作 & 學習apical 4-chamber view、血管超音波、TEE觀摩 \\
\hline
急診待命 & STEMI值班待命 & 計畫中安排但無實際急診案例 \\
\hline
\end{tabular}
\end{table}

\subsection{量化學習成果}

\begin{table}[h]
\centering
\caption{見習期間量化指標統計}
\begin{tabular}{|p{4cm}|p{2cm}|p{6cm}|}
\hline
\textbf{學習項目} & \textbf{完成數量} & \textbf{說明} \\
\hline
門診病患診療觀摩 & 70+ 例 & 涵蓋高血壓、心衰竭、冠心病、PAOD等疾病 \\
\hline
病患接診 & 6 位 & 完整執行從入院到出院的病患照護流程 \\
\hline
病歷撰寫 & 6 份 & 包括admission notes和discharge notes \\
\hline
床邊超音波執行 & 6 次 & 心臟超音波基本切面獲取與判讀 \\
\hline
動脈血管超音波 & 5 次 & 頸動脈、下肢動脈血流評估 \\
\hline
靜脈血管超音波 & 5 次 & 下肢靜脈、DVT檢查技術 \\
\hline
心導管手術觀摩 & 10+ 台 & PCI、PTA、TPM、腎上腺靜脈採樣等 \\
\hline
TEE檢查觀摩 & 2 次 & 經食道超音波及LAAO手術觀摩 \\
\hline
\end{tabular}
\end{table}

\subsection{學習亮點}

\begin{enumerate}
\item \textbf{臨床經驗豐富}：參與門診70餘例病患診療，涵蓋高血壓、心衰竭、冠心病等常見疾病
\item \textbf{實務操作充實}：獨立接診6位病患，完成病歷撰寫，累積實際臨床經驗
\item \textbf{技能操作多元}：執行16次超音波檢查（心臟6次、血管11次），建立扎實操作基礎
\item \textbf{手術觀摩多元}：觀摩包括複雜PCI、周邊血管介入、心律調節器植入等高階術式
\item \textbf{學習態度積極}：主動閱讀NEJM case records，參與self-learning activities
\item \textbf{反思能力強}：每日完整記錄學習心得與反思，展現良好的自我評估能力
\end{enumerate}

\section{臨床技能評估}

\subsection{EPA評估結果}

根據指導醫師評估，各項EPAs達成情況如下：

\begin{longtable}{|p{1cm}|p{6cm}|p{2cm}|p{5cm}|}
\caption{心臟內科EPAs評估結果} \\
\hline
\textbf{EPA} & \textbf{專業活動} & \textbf{等級} & \textbf{評估說明} \\
\hline
\endfirsthead
\hline
\textbf{EPA} & \textbf{專業活動} & \textbf{等級} & \textbf{評估說明} \\
\hline
\endhead

EPA 1 & 心臟科病史詢問與風險評估 & Level 2 & 能系統性完成病史詢問，使用適當工具 \\
\hline
EPA 2 & 心臟血管系統理學檢查 & Level 3 & 能區分收縮期與舒張期雜音，準確評估頸靜脈壓 \\
\hline  
EPA 3 & 心電圖判讀與分析 & Level 2 & 能識別常見心律不整，判讀明顯ST段變化 \\
\hline
EPA 4 & 心臟超音波基礎操作與判讀 & Level 1 & 能獲取基本標準切面，進行簡單線性測量 \\
\hline
EPA 5 & 急性冠心症診斷與初步處理 & None & 未有實際急性案例處理機會 \\
\hline
EPA 6 & 心律不整診斷與處理 & None & 未有實際心律不整處理經驗 \\
\hline
EPA 7 & 心衰竭診斷與治療 & Level 1 & 能識別心衰竭症狀徵象，知道基本藥物治療 \\
\hline
EPA 8 & 高血壓診斷與管理 & Level 2 & 能開始第一線降血壓藥物，提供基本生活型態建議 \\
\hline
EPA 9 & 心導管檢查結果判讀 & Level 2 & 能判讀標準心導管影像，評估狹窄嚴重程度 \\
\hline
EPA 10 & 心臟科病患照護協調與溝通 & Level 1 & 能進行基本病情說明，提供標準化衛教 \\
\hline
\end{longtable}

\subsection{學習成就統計}

\begin{table}[h]
\centering
\caption{EPA等級分佈統計}
\begin{tabular}{|c|c|c|}
\hline
\textbf{EPA等級} & \textbf{項目數} & \textbf{百分比} \\
\hline
Level 3 & 1 & 10\% \\
Level 2 & 4 & 40\% \\
Level 1 & 3 & 30\% \\
None & 2 & 20\% \\
\hline
\textbf{總計} & \textbf{10} & \textbf{100\%} \\
\hline
\end{tabular}
\end{table}

\section{學習品質分析}

\subsection{優勢表現}

\begin{itemize}
\item \textbf{理學檢查技能}：EPA 2達到Level 3，展現優異的心血管系統檢查能力
\item \textbf{臨床思維}：能整合多重資訊進行初步診斷判斷
\item \textbf{學習自律性}：主動學習IVUS、ECG判讀等進階技術
\item \textbf{專業態度}：與團隊建立良好溝通關係，展現專業素養
\end{itemize}

\subsection{改進建議}

\begin{itemize}
\item \textbf{急症處理經驗}：建議增加急性冠心症及心律不整的模擬訓練
\item \textbf{超音波技能}：需要更多實際操作練習以提升影像獲取能力  
\item \textbf{臨床決策}：建議參與更多複雜病例的醫療決策討論
\item \textbf{溝通技巧}：可加強困難病情告知和家屬溝通的練習
\end{itemize}

\section{標準化評分系統}

\subsection{評分架構}

本評分系統採用多元評量方式，包含三個主要評分項目：
\begin{itemize}
\item \textbf{EPA基礎能力表現 (75\%)}：基於10項EPAs的達成等級
\item \textbf{學習態度與表現 (20\%)}：包含學習積極性、專業態度、自主學習能力
\item \textbf{特殊貢獻加分 (5\%)}：創新表現、額外學習成果、卓越表現
\end{itemize}

\subsection{EPA等級評分標準}

\begin{table}[h]
\centering
\caption{EPA等級對應分數表}
\begin{tabular}{|c|c|p{8cm}|}
\hline
\textbf{EPA等級} & \textbf{對應分數} & \textbf{能力描述} \\
\hline
Level 5 & 100分 & 完全獨立執行並能指導他人 \\
\hline
Level 4 & 98分 & 可獨立執行，監督者隨時可得 \\
\hline
Level 3 & 95分 & 可獨立執行，但需回報並接受指導 \\
\hline
Level 2 & 90分 & 可在主治醫師直接監督下執行 \\
\hline
Level 1 & 85分 & 僅能觀察學習，無法獨立操作 \\
\hline
None & 75分 & 未有學習機會（見習生合理情況）\\
\hline
\end{tabular}
\end{table}

\subsection{學習態度評分準則}

\begin{table}[h]
\centering
\caption{學習態度評分標準}
\begin{tabular}{|c|p{10cm}|}
\hline
\textbf{分數區間} & \textbf{表現標準} \\
\hline
95-100分 & 學習態度極為積極，主動尋求學習機會，完整記錄學習歷程，展現優異專業素養 \\
\hline
90-94分 & 學習態度良好，能配合教學安排，有基本學習記錄，專業表現合宜 \\
\hline
85-89分 & 學習態度一般，被動參與，學習記錄不完整 \\
\hline
80-84分 & 學習態度消極，缺乏主動性，專業表現待加強 \\
\hline
\end{tabular}
\end{table}

\subsection{周煜珉醫師評分計算}

\subsubsection{EPA基礎能力分數}
\begin{align}
\text{EPA基礎分} &= \frac{(95 \times 1) + (90 \times 4) + (85 \times 3) + (75 \times 2)}{10}\\
&= \frac{95 + 360 + 255 + 150}{10}\\
&= \frac{860}{10} = 86\text{分}
\end{align}

\subsubsection{學習態度與表現評分}
根據學習記錄分析：
\begin{itemize}
\item 每日完整學習記錄：優
\item 主動參與額外學習活動：優  
\item 專業態度表現：優
\item 自我反思能力：優
\end{itemize}
\textbf{學習態度得分：98分}

\subsubsection{特殊貢獻加分}
\begin{itemize}
\item 主動閱讀NEJM case records
\item 參與self-learning echocardiography simulator
\item 學習進階技術（IVUS、DK crush技術）
\item 積極參與團隊協作
\end{itemize}
\textbf{特殊貢獻得分：100分}

\subsubsection{最終總分計算}
\begin{align}
\text{總分} &= 86 \times 0.75 + 98 \times 0.20 + 100 \times 0.05\\
&= 64.5 + 19.6 + 5.0\\
&= 89.1\text{分}
\end{align}

考量見習期間的卓越學習表現和積極態度，給予額外2.9分的指導教師裁量加分。

\textcolor{red}{\Large \textbf{最終評分：92分 (優良)}}

\subsection{評分等級對照}
\begin{table}[h]
\centering
\caption{總評成績等級表}
\begin{tabular}{|c|c|c|}
\hline
\textbf{分數區間} & \textbf{等級} & \textbf{評語} \\
\hline
90-100分 & A (優良) & 表現優異，具備良好臨床基礎 \\
\hline
80-89分 & B (良好) & 表現良好，達到預期學習目標 \\
\hline
70-79分 & C (合格) & 基本達標，仍有改進空間 \\
\hline
60-69分 & D (待改進) & 未達預期，需要額外指導 \\
\hline
60分以下 & F (不合格) & 嚴重不足，建議重新安排學習 \\
\hline
\end{tabular}
\end{table}

\section{總結與建議}

周煜珉醫師在兩週的心臟內科見習期間展現積極的學習態度和良好的專業素養。在基礎臨床技能方面表現優異，特別是理學檢查技能達到獨立執行水準。建議未來持續加強急症處理能力和臨床決策技巧的培養。

\subsection{未來發展建議}
\begin{enumerate}
\item 增加急診值班經驗，累積急性心血管疾病處理能力
\item 定期參與心導管室實習，提升介入性治療認知
\item 加強超音波操作練習，建立獨立檢查能力
\item 參與困難病例討論，提升臨床思維能力
\end{enumerate}

\vspace{1cm}
\noindent
\begin{tabular}{ll}
指導醫師簽名：& \underline{\hspace{4cm}} \\
& 謝慕揚 醫師 \\[0.5cm]
日期：& \underline{\hspace{4cm}} \\
& 2025年9月14日
\end{tabular}

\section{學生的每日學習日誌}
\includepdf[pages=-]{Clerk 周煜珉 2025-09-02 2 weeks daily note.pdf}




\end{document}